%=========== ΛΥΣΗ - 61ο ΘΕΜΑ Β ===========

\begin{thema}{Β}
Οι συναρτήσεις είναι
\[
f(x)=\ln x,\quad x>0
\]
και
\[
g(x)=a x^2+\beta,\quad x\in\mathbb{R}.
\]
Αφού οι γραφικές τους παραστάσεις έχουν κοινή εφαπτομένη στο κοινό τους σημείο με τετμημένη $x_0=1$, οι τιμές και οι παράγωγοι στο $1$ είναι ίσες.

\begin{erwthma}
\item Έχουμε
\[
f(1)=g(1)\Leftrightarrow \ln 1=a+\beta
\Leftrightarrow 0=a+\beta
\Leftrightarrow \beta=-a.
\]
Επίσης
\[
f'(x)=\frac{1}{x}
\quad\text{και}\quad
g'(x)=2ax.
\]
Άρα
\[
f'(1)=g'(1)\Leftrightarrow 1=2a
\Leftrightarrow a=\frac{1}{2}.
\]
Από τη σχέση $\beta=-a$ προκύπτει
\[
\beta=-\frac{1}{2}.
\]
Επομένως
\[
a=\frac{1}{2},\quad \beta=-\frac{1}{2}.
\]

\item Με τις τιμές του πρώτου ερωτήματος έχουμε
\[
g(x)=\frac{x^2}{2}-\frac{1}{2}.
\]
Ορίζουμε
\[
h(x)=f(x)-g(x)=\ln x-\frac{x^2}{2}+\frac{1}{2},\quad x>0.
\]
Τότε
\[
h'(x)=\frac{1}{x}-x=\frac{1-x^2}{x}
=\frac{(1-x)(1+x)}{x}.
\]
Επειδή $x>0$, έχουμε $x>0$ και $1+x>0$. Άρα το πρόσημο της $h'(x)$ καθορίζεται από το $1-x$:
\begin{itemize}
\item $h'(x)>0\Leftrightarrow 0<x<1$,
\item $h'(x)=0\Leftrightarrow x=1$,
\item $h'(x)<0\Leftrightarrow x>1$.
\end{itemize}
Επομένως η $h$ είναι \textbf{γνησίως αύξουσα} στο $(0,1]$ και \textbf{γνησίως φθίνουσα} στο $[1,+\infty)$.

Στο $x=1$ παρουσιάζει ολικό μέγιστο με τιμή
\[
h(1)=\ln 1-\frac{1}{2}+\frac{1}{2}=0.
\]

\item Από το προηγούμενο ερώτημα η μέγιστη τιμή της $h$ είναι $0$. Άρα
\[
h(x)\leq 0,\quad x>0.
\]
Δηλαδή
\[
\ln x-\frac{x^2}{2}+\frac{1}{2}\leq 0
\Leftrightarrow
\ln x\leq \frac{x^2-1}{2}.
\]
Επομένως
\[
\ln x\leq \frac{x^2-1}{2},\quad x>0.
\]

\item Στο διάστημα $[1,e]$ ισχύει $h(x)\leq 0$. Άρα το ζητούμενο εμβαδόν είναι
\[
E=\dint_1^e |h(x)|\d x=\dint_1^e (-h(x))\d x.
\]
Επομένως
\begin{align*}
E
&=\dint_1^e\left(\frac{x^2}{2}-\frac{1}{2}-\ln x\right)\d x\\
&=\left[\frac{x^3}{6}-\frac{x}{2}-x\ln x+x\right]_1^e.
\end{align*}
Υπολογίζουμε:
\[
\left[\frac{x^3}{6}-\frac{x}{2}-x\ln x+x\right]_1^e
=\left(\frac{e^3}{6}-\frac{e}{2}-e+e\right)
-\left(\frac{1}{6}-\frac{1}{2}-0+1\right).
\]
Άρα
\[
E=\frac{e^3}{6}-\frac{e}{2}-\frac{2}{3}
=\frac{e^3-3e-4}{6}.
\]
\end{erwthma}
\end{thema}

%=========== ΤΕΛΟΣ ΛΥΣΗΣ - 61ο ΘΕΜΑ Β ===========
